\documentclass{amsart}
\usepackage[margin=1in]{geometry}
\usepackage{tikz,tikz-cd}
\usetikzlibrary{calc}
\usepackage[colorlinks=true]{hyperref}
\usepackage{ifthen,calc,amsmath,amssymb,mathtools,enumitem}
\usepackage{xcolor,amssymb}

\newtheorem{lemma}{Lemma}
\newtheorem{proposition}{Proposition}
\newtheorem{corollary}{Corollary}
\newtheorem*{problem*}{Problem}
\theoremstyle{definition}
\newtheorem{remark}{Remark}
\newtheorem{definition}{Definition}

\newcommand{\from}{\colon}
\newcommand{\RR}{\mathbb{R}}
\newcommand{\ZZ}{\mathbb{Z}}
\newcommand{\defeq}{\coloneqq}
\newcommand{\bdy}{\partial}

\begin{document}

\begin{problem*}
  Let $G=\langle a_1,\dots,a_r\rangle$ be a free group, $L=\{a_1^{\pm1},\dots,a_r^{\pm1}\}$.
  Fix a reducible quasi-reduced word $w=w_1\cdots w_{2n}$.
  A matching $M$ consists of $n$ properly embedded intervals in $\RR\times\RR_+$ with
  $\bdy M=\{1,\dots,2n\}\times 0$, pairwise transverse with no triple points, such that:
  for each component $R$ of $\RR\times\RR_+\setminus M$ and each component $A$ of
  $M\cap\bdy\overline R$, the boundary $\bdy A$ consists of two points $(i,0),(j,0)$ with
  $w_i=w_j^{-1}$.
  Let $F_s$ be the CW complex whose $k$-cells are isotopy classes of $k$-crossing matchings,
  with the cube structure and resolution attaching maps described in the problem.
  Is $F_s$ contractible?
\end{problem*}

\begin{proof}[Solution]
The answer is \textbf{no}: $F_s$ is not contractible (indeed not $2$-connected) in general.
We exhibit a word $w$ for which the complex $F_w$ is homeomorphic to the $2$-sphere $S^2$.

Throughout, take $r\ge 3$, write $a=a_1,\ b=a_2,\ c=a_3$, and fix
\[
  w \;=\; a\,a^{-1}\,b\,b^{-1}\,a\,a^{-1}\,c\,c^{-1}\,a\,a^{-1},
  \qquad n=5 .
\]
This word is reducible ($w=1$ in $G$) and quasi-reduced (it never contains a substring
$\ell\ell^{-1}\ell$). The letters at the ten positions are
\[
  w_1=a,\ w_2=a^{-1},\ w_3=b,\ w_4=b^{-1},\ w_5=a,\ w_6=a^{-1},
  \ w_7=c,\ w_8=c^{-1},\ w_9=a,\ w_{10}=a^{-1}.
\]
The inverse relation $w_i=w_j^{-1}$ holds precisely in the following cases:
\begin{equation}\label{eq:inv}
\begin{aligned}
&\text{one of }i,j\text{ is an }a\text{-position }\{1,5,9\}\text{ and the other an }
 a^{-1}\text{-position }\{2,6,10\};\\
&\text{or }\{i,j\}=\{3,4\};\qquad\text{or }\{i,j\}=\{7,8\}.
\end{aligned}
\end{equation}

We organize the proof as follows.
Section~\ref{sec:bigon} gives two structural lemmas that reduce all admissible matchings to a
finite, geometrically minimal list.
Section~\ref{sec:enum} carries out the enumeration: there are exactly five $0$-cells,
six $1$-cells, three $2$-cells, and no cells of dimension $\ge 3$.
Section~\ref{sec:attach} computes the attaching maps, and Section~\ref{sec:sphere} identifies
$F_w$ with $S^2$.

\subsection{Two reduction lemmas}\label{sec:bigon}

An \emph{arc} of $M$ is one of its $n=5$ properly embedded intervals; its two boundary points
are positions $i<j$, its \emph{ends}. We call the third bullet of the matching axioms the
\emph{region condition}: for each region $R$ and each component $A$ of $M\cap\bdy\overline R$,
the two endpoints of $A$ are positions $i,j$ with $w_i=w_j^{-1}$.

\begin{lemma}[No closed sub-loops; no bigons]\label{lem:bigon}
Let $M$ be a matching satisfying the region condition. Then:
\begin{enumerate}
\item[(i)] For every region $R$ and every component $A$ of $M\cap\bdy\overline R$, both
endpoints of $A$ lie on $\RR\times0$. In particular no such $A$ is a closed loop and no $A$
has a crossing point as an endpoint.
\item[(ii)] No two arcs of $M$ cross each other more than once.
\end{enumerate}
\end{lemma}

\begin{proof}
(i) A component $A$ of $M\cap\bdy\overline R$ is a connected sub-segment of a single arc,
maximal among those lying on $\bdy\overline R$. Each of its two endpoints is a point of $M$
that is extreme on $A$, hence is either an end of the arc (a point of
$\bdy M=\{1,\dots,2n\}\times0$) or a transverse crossing point. The region condition asserts
$\bdy A$ consists of two points $(i,0),(j,0)$ on $\RR\times0$. Hence neither endpoint is an
interior crossing point, and $A$ is not a closed loop (which has empty boundary).

(ii) Suppose two arcs $\alpha,\beta$ cross at least twice. Among all sub-bigons cut out by
$\alpha\cup\beta$, choose one, $B$, that is \emph{innermost}, i.e.\ whose interior contains no
crossing point of $M$; such a bigon exists because any crossing inside a bigon subdivides it
into smaller bigons, and we descend to a minimal one. Then $\operatorname{int}B$ meets no arc
of $M$, so $B$ is the closure of a region $R$, and $\bdy\overline R$ is the union of a
sub-arc of $\alpha$ and a sub-arc of $\beta$, meeting at the two crossing points $p,q$. Thus a
component $A$ of $M\cap\bdy\overline R$ has $\bdy A\subseteq\{p,q\}$, with $p,q$ interior
crossing points not on $\RR\times0$, contradicting~(i).
\end{proof}

By Lemma~\ref{lem:bigon}(ii), in any admissible $M$ each pair of arcs crosses at most once, so
the crossing number of $M$ equals the number of \emph{interleaving} pairs of arcs (pairs
$\{i_1,j_1\},\{i_2,j_2\}$ with $i_1<i_2<j_1<j_2$). Two embedded arcs whose ends interleave
cross exactly once, and a system of arcs with prescribed ends and exactly one crossing per
interleaving pair is unique up to isotopy. Hence the isotopy class of an admissible matching
is determined by its pairing of $\{1,\dots,10\}$ alone, and we may identify admissible cells
with admissible pairings, the dimension being the interleaving count.

\begin{lemma}[Forced short arcs]\label{lem:forced}
In every admissible matching $M$, position $3$ is matched to $4$ and position $7$ to $8$.
\end{lemma}

\begin{proof}
Position $3$ carries the letter $b$, and by~\eqref{eq:inv} the only $j$ with $w_3=w_j^{-1}$ is
$j=4$. Let $\alpha$ be the arc through position $3$, let $R$ be a region adjacent to $\alpha$
near the end $(3,0)$, and let $A\subseteq\alpha$ be the component of $M\cap\bdy\overline R$
containing $(3,0)$. By Lemma~\ref{lem:bigon}(i) the other endpoint of $A$ is a position
$(j,0)$ with $w_3=w_j^{-1}$, forcing $j=4$; thus $A$ runs from $(3,0)$ to $(4,0)$ along
$\alpha$. If $A\subsetneq\alpha$ then $\alpha$ would pass through $(4,0)$ as an interior
point, but $(4,0)\in\bdy M$ is an arc end, not interior. Hence $A=\alpha$ and $\alpha$ joins
$3$ to $4$. The same argument at position $7$ (letter $c$, inverse only at $8$) forces the arc
through $7$ to join $7$ to $8$.
\end{proof}

By Lemma~\ref{lem:forced}, two arcs are always $\{3,4\}$ and $\{7,8\}$; being between adjacent
positions, they cross no other arc. The remaining three arcs match the six positions
\[
  P\defeq\{1,2,5,6,9,10\},\qquad a\text{-positions }\{1,5,9\},\ \ a^{-1}\text{-positions }\{2,6,10\}.
\]

\subsection{Well-definedness of $F_s$ as a cubical CW complex}\label{sec:welldef}

Before extracting cell counts we verify that the data of the Problem assemble into a genuine
CW (indeed cubical) complex. We prove the three assertions implicit in the construction:
\begin{enumerate}
\item[(W1)] resolving a coherent $k$-crossing matching $M$ at a crossing $c$, in either local
manner $a\in\bdy I_c$, gives a matching $M_a$ with crossing set $X(M_a)=X(M)\setminus\{c\}$;
\item[(W2)] $M_a$ is again \emph{coherent} (it satisfies the region condition);
\item[(W3)] the facet--to--cell assignment is well defined, consistent under iterated resolution,
and continuous, so the cubes glue to a CW complex.
\end{enumerate}

Throughout, ``matching'' means $n$ properly embedded intervals in $H\defeq\RR\times\RR_+$ with
$\bdy M=\{1,\dots,2n\}\times0$, pairwise transverse, no triple points; $M$ is \emph{coherent} if
it satisfies the region condition (third matching axiom), and $X(M)$ is its set of transverse
self–intersections. We use Lemma~\ref{lem:bigon} (proved from the region condition only): in a
coherent matching no two arcs cross more than once and no region of $H\setminus M$ is bounded
entirely by arcs of $M$ — equivalently every boundary component of a region runs between two
points of $\RR\times0$.

\subsubsection{Anatomy of a boundary component}\label{ssec:anatomy}

Fix a coherent $M$ and a region $R$ of $H\setminus M$ with closure $\overline R$. Let $A$ be a
connected component of $M\cap\bdy\overline R$.

\begin{lemma}[Shape of $A$]\label{lem:shapeA}
$A$ is an embedded arc; its interior is a union of arc–segments of $M$ meeting only at
transverse crossing points, and at each such crossing $A$ ``turns the corner'' (passes from one
crossing strand to the other). Both endpoints of $A$ lie on $\RR\times0$ and are arc–ends
$(i,0),(j,0)$ with $w_i=w_j^{-1}$.
\end{lemma}

\begin{proof}
The frontier $\bdy\overline R$ of a planar region lies in the walls $M\cup(\RR\times0)$, so it is
a piecewise–smooth $1$–manifold and $M\cap\bdy\overline R$ is a disjoint union of compact arcs
and circles; the $A$'s are its components. By Lemma~\ref{lem:bigon} none is a circle and each has
its two boundary points on $\RR\times0$; the region condition gives $w_i=w_j^{-1}$ at those ends.

For the interior: let $x$ be an interior point of $A$, so $x\in M\cap\bdy\overline R$,
$x\notin\RR\times0$. If $x$ is a smooth point of $M$, then $M$ is locally one strand splitting a
disk into two half–disks, exactly one inside $R$, so $\bdy\overline R\cap M$ is that single
strand near $x$ and $A$ continues smoothly. If $x$ is a crossing of strands $\sigma,\tau$, the
four local sectors are traces of (at most) four regions; $R$ occupies one sector, whose two
bounding rays lie on $\sigma$ and $\tau$. Hence near $x$, $\bdy\overline R\cap M$ is the corner
of one ray of $\sigma$ and one ray of $\tau$: $A$ enters along $\sigma$ and leaves along $\tau$,
turning the corner, and $x$ is interior to $A$.
\end{proof}

\begin{definition}\label{def:corners}
Let $c$ be a crossing of strands $\alpha,\beta$ of $M$ and $D=D_c\subset H$ a closed round disk
so small that $D\cap M=\alpha_0\cup\beta_0$ is two chords crossing once at $c$, with
$\bdy D\cap M=\{\alpha^+,\beta^+,\alpha^-,\beta^-\}$ in this cyclic order (the arcs alternate,
by transversality). The four sectors of $D\setminus(\alpha_0\cup\beta_0)$, cyclically, are
$S_{++}$ (between $\alpha^+,\beta^+$), $S_{-+}$ (between $\beta^+,\alpha^-$),
$S_{--}$ (between $\alpha^-,\beta^-$), $S_{+-}$ (between $\beta^-,\alpha^+$); each lies in a
region of $M$, with opposite sectors $S_{++},S_{--}$ in regions $R_1,R_3$ and $S_{-+},S_{+-}$ in
regions $R_2,R_4$ (a priori not all distinct).
\end{definition}

\subsubsection{(W1): locality of crossings}\label{ssec:W1}

Resolving $M$ at $c$ deletes $\alpha_0\cup\beta_0$ and inserts one of the two pairs of disjoint
chords of $D$ with the same endpoints on $\bdy D$:
\[
  a{=}0:\ \gamma=\overline{\alpha^+\beta^+},\ \delta=\overline{\alpha^-\beta^-};
  \qquad
  a{=}1:\ \gamma'=\overline{\beta^+\alpha^-},\ \delta'=\overline{\alpha^+\beta^-}.
\]
Outside $D$ nothing changes. Since the endpoints on $\bdy D$ lie on the unchanged part of $M$,
the result $M_a$ is again $n$ properly embedded intervals with the same boundary
$\{1,\dots,2n\}\times0$; a small isotopy keeps it transverse with no triple point.

\begin{lemma}[(W1)]\label{lem:W1}
$X(M_a)=X(M)\setminus\{c\}$ as subsets of $H$, for either $a$.
\end{lemma}

\begin{proof}
$M$ and $M_a$ agree outside $\operatorname{int}D$, where they have the same crossings
$X(M)\setminus\{c\}$ (the disk contains only $c$). Inside $D$, $M$ has the single crossing $c$
while $M_a\cap D$ is a pair of disjoint chords with no crossing. Hence
$X(M_a)=X(M)\setminus\{c\}$.
\end{proof}

This is purely geometric and needs no coherence: smoothing one crossing removes that crossing
and disturbs no other. The real content is the coherence statement (W2).

\subsubsection{(W2): resolution preserves coherence}\label{ssec:W2}

Take $a=0$; the case $a=1$ is identical with the two diagonal pairs exchanged. Inside $D$ the
smoothing reconnects germs as $\gamma\colon\alpha^+\!-\!\beta^+$ and
$\delta\colon\alpha^-\!-\!\beta^-$.

\paragraph{Local picture of regions.}
The smoothing changes only the regions occupying a sector at $c$. Replacing the cross
$\alpha_0\cup\beta_0$ by $\gamma\cup\delta$ has the standard planar effect:
\begin{itemize}
\item sectors $S_{-+}$ and $S_{+-}$ become joined through the open channel of $D$ between
$\gamma$ and $\delta$, so regions $R_2,R_4$ are \emph{merged} into one region $R_{24}$ of $M_a$;
\item sectors $S_{++},S_{--}$ stay separate (each now cut off by one chord), so $R_1,R_3$
persist as regions $R_1',R_3'$ of $M_a$, reshaped only inside $D$.
\end{itemize}
(The smoothing $a=1$ instead merges $R_1,R_3$ and keeps $R_2,R_4$ separate.) No other region is
created or destroyed; outside $D$ all regions abut $\bdy D$ along the same arcs as before. Thus
the regions of $M_a$ are: the regions of $M$ disjoint from $\operatorname{int}D$ (unchanged),
the persisting $R_1',R_3'$, and the merged $R_{24}$.

\paragraph{Boundary components near $c$.}
We track boundary components directly; we do \emph{not} claim $M_a$ is isotopic to $M$ (a
smoothing genuinely changes the $1$--manifold inside $D$). All boundary components of regions of
$M$ disjoint from $\operatorname{int}D$ are also boundary components of $M_a$, and conversely,
since $M=M_a$ outside $\operatorname{int}D$; only components that meet $D$ are affected.

By Lemma~\ref{lem:shapeA}, near the crossing $c$ a boundary component of a region of $M$ either
avoids $c$ or passes once through $c$, turning the corner. List the at most four boundary
components of $M$ that pass through $c$, one per sector:
$A_1$ (sector $S_{++}$, on $R_1$), $A_2$ ($S_{-+}$, on $R_2$), $A_3$ ($S_{--}$, on $R_3$),
$A_4$ ($S_{+-}$, on $R_4$). At the corner $c$ each $A_t$ connects the two strand--germs bounding
its sector:
\[
  A_1:\ \alpha^+\!\leftrightarrow\!\beta^+,\quad
  A_2:\ \beta^+\!\leftrightarrow\!\alpha^-,\quad
  A_3:\ \alpha^-\!\leftrightarrow\!\beta^-,\quad
  A_4:\ \beta^-\!\leftrightarrow\!\alpha^+ .
\]
For a germ $g\in\{\alpha^+,\beta^+,\alpha^-,\beta^-\}$, denote by $E(g)\in\RR\times0$ the
$x$--axis end reached by following the strand of $M$ outward from $g$, away from $c$ along
$\bdy\overline R$, until it hits the $x$--axis; this exists and is well defined by
Lemma~\ref{lem:shapeA} (the boundary component continues until an $x$--axis end). The corner
component $A_t$ joins the two ends of its germ pair; e.g.\ $A_1$ joins $E(\alpha^+)$ to
$E(\beta^+)$, so by the region condition $w_{E(\alpha^+)}=w_{E(\beta^+)}^{-1}$, and similarly for
$A_2,A_3,A_4$. (If some sector carries no boundary component through $c$, the corresponding $A_t$
is absent; the argument below uses only those relations that are present, together with the
type--colouring, and is unaffected.)

\paragraph{Effect of the smoothing on these components.}
After the $a{=}0$ smoothing the corner at $c$ is gone; inside $D$ the boundary now follows the
two disjoint chords $\gamma\ (\alpha^+\!-\!\beta^+)$ and $\delta\ (\alpha^-\!-\!\beta^-)$, and $D$
contains no crossing of $M_a$. A boundary component of $M_a$ that meets $D$ agrees with $M$ outside
$\operatorname{int}D$ and, inside $D$, runs along a single chord ($\gamma$ or $\delta$), connecting
that chord's two germs. Thus the half--arcs of $M$ leading out of $D$ from the four germs are
re--spliced by $M_a$ along the pairs $\{\alpha^+,\beta^+\}$ (via $\gamma$) and
$\{\alpha^-,\beta^-\}$ (via $\delta$). On the persisting side, $\gamma$ recreates $A_1$ (same ends
$E(\alpha^+),E(\beta^+)$) and $\delta$ recreates $A_3$; these are the through--$D$ boundary
components of $R_1',R_3'$. The merged side is analysed in the proof of Lemma~\ref{lem:splice},
where we determine the entire frontier of $R_{24}$.

\begin{lemma}[Splicing preserves inverse ends]\label{lem:splice}
Let $A'$ be a component of $M_a\cap\bdy\overline R$ for any region $R$ of $M_a$, with ends
$(i,0),(j,0)$ on $\RR\times0$. Then $w_i=w_j^{-1}$, and $A'$ is an embedded arc (no circle
component arises).

\end{lemma}

\begin{proof}
If $R$ is disjoint from $\operatorname{int}D$, then $A'$ is a boundary component of a region of
$M$ and the claim is the region condition for $M$. So let $R\in\{R_1',R_3',R_{24}\}$ and use the
germ bookkeeping above. Recall $E(g)$ denotes the $x$--axis end reached by following $M$ (equally
$M_a$, since they agree outside $D$) out of $D$ from germ $g$, and that coherence of $M$ gives,
for the four corner components $A_t$ of $M$,
\begin{equation}\label{eq:germinv}
  w_{E(\alpha^+)}=w_{E(\beta^+)}^{-1},\qquad
  w_{E(\beta^+)}=w_{E(\alpha^-)}^{-1},\qquad
  w_{E(\alpha^-)}=w_{E(\beta^-)}^{-1},\qquad
  w_{E(\beta^-)}=w_{E(\alpha^+)}^{-1},
\end{equation}
these being the inverse--end conditions for $A_1,A_2,A_3,A_4$ respectively (whenever the
corresponding corner component exists). In particular all four ends $E(\alpha^\pm),E(\beta^\pm)$
represent the \emph{same} letter up to sign: writing $\ell\defeq w_{E(\alpha^+)}$, the four
relations give $w_{E(\beta^+)}=\ell^{-1}$, $w_{E(\alpha^-)}=\ell$, $w_{E(\beta^-)}=\ell^{-1}$.

\emph{Persisting region.} A component $A'$ of $R_1'$ through $D$ runs along $\gamma$ and so joins
$E(\alpha^+)$ to $E(\beta^+)$; by~\eqref{eq:germinv} (the $A_1$ relation) its ends are inverse.
For $R_3'$, $A'$ runs along $\delta$, joining $E(\alpha^-)$ to $E(\beta^-)$, inverse by the $A_3$
relation. Components of $R_1',R_3'$ not meeting $D$ are unchanged components of $M$, inverse by
coherence. Each is an arc, being a re--splice of arcs along a single chord.

\emph{Merged region.} We pin down the frontier of $R_{24}$ exactly. Colour the germs by
\emph{type}: $\alpha^+,\alpha^-$ have type $\mathsf A$ and $\beta^+,\beta^-$ type $\mathsf B$. By
the computation above, $w_{E(g)}=\ell$ if $\operatorname{type}(g)=\mathsf A$ and $w_{E(g)}=\ell^{-1}$
if $\operatorname{type}(g)=\mathsf B$. Hence \emph{a boundary component has inverse ends as soon as
its two terminal germ--ends have opposite type.}

Now describe $\bdy\overline{R_{24}}$. Inside $D$, $R_{24}$ (the channel) is bounded by the
$R_{24}$--side of $\gamma$, joining germs $\alpha^+,\beta^+$, and the $R_{24}$--side of $\delta$,
joining $\alpha^-,\beta^-$; inside $D$ there are no crossings of $M_a$, so any boundary arc of
$R_{24}$ that enters $D$ runs along exactly one of these two chords and exits at its partner germ.
Outside $D$, the frontier of $R_{24}$ is the old frontier of $R_2$ (which meets $\bdy D$ at germs
$\beta^+,\alpha^-$) and of $R_4$ (which meets $\bdy D$ at germs $\beta^-,\alpha^+$). Matching
in--$D$ and out--of--$D$ pieces at the shared germs ($\alpha^+$: $\gamma\leftrightarrow R_4$;
$\beta^+$: $\gamma\leftrightarrow R_2$; $\alpha^-$: $\delta\leftrightarrow R_2$; $\beta^-$:
$\delta\leftrightarrow R_4$), the closed frontier of $R_{24}$ is the single cyclic walk
\[
  \alpha^+ \xrightarrow{\ \gamma\ } \beta^+ \xrightarrow{\ \bdy R_2\ } \alpha^-
          \xrightarrow{\ \delta\ } \beta^- \xrightarrow{\ \bdy R_4\ } \alpha^+ ,
\]
in which consecutive germs are joined by a chord or by an external frontier piece. A maximal
arc--portion $A'$ of $M_a$ on $\bdy\overline{R_{24}}$ is a segment of this walk between two
consecutive $x$--axis touches; its two terminal germ--ends are two germs that are \emph{adjacent}
in the displayed $4$--cycle (the $x$--axis touches lie in the interior of the external pieces and
do not change which two germs flank a maximal arc run). But in the $4$--cycle
$\alpha^+,\beta^+,\alpha^-,\beta^-$ the types alternate $\mathsf A,\mathsf B,\mathsf A,\mathsf B$,
so any two adjacent germs have opposite type. Hence the two ends of $A'$ have opposite type, and
$w_i=w_j^{-1}$. Finally each external piece reaches the $x$--axis (Lemma~\ref{lem:shapeA} for
$R_2,R_4$), so no boundary component of $R_{24}$ is an all--arc circle: every $A'$ has two genuine
$x$--axis endpoints.

(The complementary smoothing $a{=}1$ is identical after exchanging the diagonal pairs
$\{S_{++},S_{--}\}\leftrightarrow\{S_{-+},S_{+-}\}$, i.e.\ exchanging the roles of the relations
$A_1,A_3$ with $A_2,A_4$ in~\eqref{eq:germinv}.)
\end{proof}


\begin{proposition}[(W2)]\label{prop:W2}
If $M$ is coherent then $M_a$ is coherent for both $a\in\{0,1\}$.
\end{proposition}

\begin{proof}
By Lemma~\ref{lem:splice}, every boundary component of every region of $M_a$ is an arc with both
ends on $\RR\times0$ and with $w_i=w_j^{-1}$, and no region of $M_a$ has an all–arc (circle)
boundary component. That is exactly the region condition: $M_a$ is coherent.
\end{proof}

\begin{corollary}[Dimension drop]\label{cor:dimdrop}
For coherent $M$ and either resolution, $|X(M_a)|=|X(M)|-1$, and the cube
$C(M_a)=\prod_{d\in X(M_a)}I_d$ has dimension $\dim C(M)-1$. Equivalently, the interleaving count
of the pairing of $M_a$ is one less than that of $M$.
\end{corollary}

\begin{proof}
$|X(M_a)|=|X(M)|-1$ by Lemma~\ref{lem:W1}. By Proposition~\ref{prop:W2} $M_a$ is coherent, so by
Lemma~\ref{lem:bigon} it is bigon–free; thus its crossing number equals its interleaving count
(each interleaving pair crosses exactly once, each crossing comes from an interleaving pair), and
$\dim C(M_a)=|X(M_a)|=|X(M)|-1=\dim C(M)-1$.

(Combinatorially this also explains why the naive ``smooth a crossing'' operation can only drop
the interleaving count by one for coherent inputs: a drop by more than one would require three
arcs of $M$ pairwise interleaving through $c$, but such a triple bounds a triangular region with
all–arc boundary, contradicting Lemma~\ref{lem:bigon}.)
\end{proof}

\subsubsection{(W3): consistency of the face maps}\label{ssec:W3}

Fix a coherent $k$–crossing matching $M$, cube $C(M)=\prod_{c\in X(M)}I_c$, each $I_c=[0,1]$ with
endpoints labelled by the two resolutions $a\in\{0,1\}$. Its codimension–one faces are
$F_{c,a}=\{a\}\times\prod_{d\in X(M)\setminus\{c\}}I_d$, for $c\in X(M)$, $a\in\{0,1\}$.

\begin{lemma}[Faces map to cells]\label{lem:facemap}
For each $(c,a)$ there is a canonical homeomorphism
$\Phi_{c,a}\colon F_{c,a}\xrightarrow{\cong}C(M_{c,a})$, where $M_{c,a}$ is the resolution of $M$
at $c$ in manner $a$, given by the identity on coordinates under the equality of index sets
$X(M)\setminus\{c\}=X(M_{c,a})$ from Lemma~\ref{lem:W1}.
\end{lemma}

\begin{proof}
By Lemma~\ref{lem:W1} the smoothing inside $D_c$ leaves every other crossing $d\neq c$ — and its
two local resolutions, hence the labelled interval $I_d$ — unchanged. So
$F_{c,a}=\prod_{d\neq c}I_d$ and $C(M_{c,a})=\prod_{d\in X(M_{c,a})}I_d$ are the same product of
the same labelled intervals, and $\Phi_{c,a}$ is the identity on coordinates, a homeomorphism. By
Proposition~\ref{prop:W2}, $M_{c,a}$ is a coherent $(k-1)$–crossing matching, so $C(M_{c,a})$ is
a $(k-1)$–cell of $F_s$.
\end{proof}

\begin{lemma}[Commuting resolutions]\label{lem:commute}
For distinct crossings $c\neq c'$ and $a,a'\in\{0,1\}$,
$(M_{c,a})_{c',a'}=(M_{c',a'})_{c,a}$ up to isotopy, with crossing set
$X(M)\setminus\{c,c'\}$, and both iterated resolutions are coherent.
\end{lemma}

\begin{proof}
Smoothings at $c$ and $c'$ are supported in disks $D_c,D_{c'}$, which we may take disjoint
(distinct points, no triple points). Modifications in disjoint disks commute, so the two orders
give the same matching, namely $M$ with both crossings smoothed; its crossing set is
$X(M)\setminus\{c,c'\}$ by two uses of Lemma~\ref{lem:W1}, and it is coherent by two uses of
Proposition~\ref{prop:W2}. The coordinate identifications are the identity on
$X(M)\setminus\{c,c'\}$ in either order.
\end{proof}

\begin{proposition}[(W3)]\label{prop:W3}
The cells $C(M)$, one per isotopy class of coherent matching, with attaching maps given on facets
by $\Phi_{c,a}$, form a regular CW complex $F_s$ of cubes in which the boundary of each $k$–cube
is glued homeomorphically onto a subcomplex of the $(k-1)$–skeleton.
\end{proposition}

\begin{proof}
Build $F_s$ by skeleta. The $0$–cells are the coherent $0$–crossing matchings (finitely many
isotopy classes, as a coherent matching is bigon–free hence determined by its pairing,
Lemma~\ref{lem:bigon}). Assume $F_s^{(k-1)}$ built from the cubes $C(M')$ with $|X(M')|\le k-1$.
For a coherent $k$–crossing matching $M$, define $\varphi_M\colon\bdy C(M)\to F_s^{(k-1)}$ on each
facet $F_{c,a}$ as $\Phi_{c,a}$ followed by the characteristic map of $C(M_{c,a})$.

Each $\Phi_{c,a}$ is a homeomorphism of the facet onto a $(k-1)$–cube (Lemma~\ref{lem:facemap}),
hence continuous. The facets of $C(M)$ overlap along codimension–two faces $F_{c,a}\cap F_{c',a'}$;
there the two prescriptions agree, since by Lemma~\ref{lem:commute} both restrict to the identity
coordinate map onto the same cell $C\big((M_{c,a})_{c',a'}\big)$, which sits in $F_s^{(k-1)}$ as
the common facet of $C(M_{c,a})$ and $C(M_{c',a'})$. Hence the facet maps glue to a single
continuous map $\varphi_M$. Iterating Lemma~\ref{lem:commute}, $\varphi_M$ maps each face of
$C(M)$ homeomorphically onto a cell of $F_s^{(k-1)}$, so it is a homeomorphism of $\bdy C(M)$ onto
a subcomplex, injective on the open cube — the condition for attaching a regular $k$–cell. Set
$F_s^{(k)}=F_s^{(k-1)}\cup_{\varphi_M}\bigsqcup_M C(M)$ over the coherent $k$–crossing classes.

The result $F_s=\bigcup_k F_s^{(k)}$ is a regular CW complex of cubes whose attaching maps are the
resolution face maps; it is closure–finite with the weak topology by construction, and finite in
each degree because coherent matchings are determined by finitely many pairings.
\end{proof}

\begin{remark}
Propositions~\ref{prop:W2} and~\ref{prop:W3} make the cell counts of the Solution well posed: the
$k$–cells of $F_s$ are exactly the isotopy classes of coherent $k$–crossing matchings, and facet
incidences are planar smoothings. For $w=aa^{-1}bb^{-1}aa^{-1}cc^{-1}aa^{-1}$ the attaching data
of Section~\ref{sec:attach} are precisely the maps $\Phi_{c,a}$ of Lemma~\ref{lem:facemap}.
\end{remark}


\subsection{Enumeration of cells}\label{sec:enum}

By Lemmas~\ref{lem:bigon} and~\ref{lem:forced}, an admissible matching is an embedded
representative of a perfect matching of $P=\{1,2,5,6,9,10\}$ into three arcs, together with the
two fixed arcs $\{3,4\},\{7,8\}$; its dimension (crossing number) equals the number of
interleaving pairs among the three arcs on $P$. There are $5!!=15$ perfect matchings of six
points. We will prove that exactly $14$ of them are admissible, distributed as follows (we
suppress the fixed arcs and list a matching by its three arcs on $P$):
\begin{equation}\label{eq:list}
\renewcommand{\arraystretch}{1.25}
\begin{array}{l@{\quad}l}
\textbf{$0$ crossings (five $0$-cells):}
 & Q_1=\{1,2\}\{5,6\}\{9,10\},\ \ Q_2=\{1,2\}\{5,10\}\{6,9\},\\
 & Q_3=\{1,6\}\{2,5\}\{9,10\},\ \ Q_4=\{1,10\}\{2,5\}\{6,9\},\\
 & Q_5=\{1,10\}\{2,9\}\{5,6\};\\[2pt]
\textbf{$1$ crossing (six $1$-cells):}
 & E_1=\{1,9\}\{2,5\}\{6,10\},\ \ E_2=\{1,5\}\{2,6\}\{9,10\},\\
 & E_3=\{1,2\}\{5,9\}\{6,10\},\ \ E_4=\{1,10\}\{2,6\}\{5,9\},\\
 & E_5=\{1,9\}\{2,10\}\{5,6\},\ \ E_6=\{1,5\}\{2,10\}\{6,9\};\\[2pt]
\textbf{$2$ crossings (three $2$-cells):}
 & T_1=\{1,6\}\{2,10\}\{5,9\},\ \ T_2=\{1,9\}\{2,6\}\{5,10\},\\
 & T_3=\{1,5\}\{2,9\}\{6,10\};\\[2pt]
\textbf{inadmissible:} & N=\{1,6\}\{2,9\}\{5,10\}.
\end{array}
\end{equation}

\paragraph{Crossing counts.}
The five matchings $Q_1,\dots,Q_5$ are exactly the $C_3=5$ non-crossing (Catalan) matchings of
six points; none has an interleaving pair. Each $E_k$ has exactly one interleaving pair, each
$T_m$ exactly two, and $N$ has all three pairs interleaving; these are the only possibilities,
since three arcs have at most $\binom32=3$ pairs. The interleaving pairs are, explicitly,
\[
\begin{array}{lll}
E_1:\{1,9\}\!\times\!\{6,10\}, & E_2:\{1,5\}\!\times\!\{2,6\}, & E_3:\{5,9\}\!\times\!\{6,10\},\\
E_4:\{2,6\}\!\times\!\{5,9\}, & E_5:\{1,9\}\!\times\!\{2,10\}, & E_6:\{1,5\}\!\times\!\{2,10\},\\
T_1:\{1,6\}\!\times\!\{2,10\},\{1,6\}\!\times\!\{5,9\}, & T_2:\{1,9\}\!\times\!\{5,10\},\{2,6\}\!\times\!\{5,10\}, & T_3:\{1,5\}\!\times\!\{2,9\},\{2,9\}\!\times\!\{6,10\},\\
N:\{1,6\}\!\times\!\{2,9\},\{1,6\}\!\times\!\{5,10\},\{2,9\}\!\times\!\{5,10\}, & &
\end{array}
\]
each verified by the defining inequality $e_0<e_1<e_2<e_3$. (For instance in $E_1$, the pair
$\{1,9\},\{6,10\}$ interleaves because $1<6<9<10$, while $\{2,5\}\subset(1,9)$ is nested, hence
non-interleaving.) Thus the dimension of the cell of each listed matching is as stated, and
\emph{no} matching of $P$ has interleaving count $\ge 3$ except $N$ (which we now exclude).

\paragraph{A parity lemma.}
The heart of the count is the following enclosure principle, which we prove for an arbitrary
matching (coherence is \emph{not} assumed).

\begin{lemma}[Even enclosure]\label{lem:even}
Let $M$ be any matching (in the sense of \S\ref{sec:welldef}: $n$ properly embedded, pairwise
transverse arcs in $H=\RR\times\RR_+$, ends $\{1,\dots,2n\}\times0$, no triple points). Let
$A\subseteq M$ be an embedded arc whose two endpoints are axis points $(i,0),(j,0)$ with $i<j$
and whose interior lies in the open half-plane and contains no arc-end of $M$. Then the number
of arc-ends of $M$ lying in the open interval $(i,j)$ is \emph{even}.
\end{lemma}

\begin{proof}
The union $\Gamma=A\cup\big([i,j]\times0\big)$ is a Jordan curve in the closed half-plane,
bounding a closed disk $D\subset H$ whose interior lies in the open half-plane; let
$U=\operatorname{int}D$. We count, modulo $2$, the arc-ends in $(i,j)$ in two ways.

Let $\gamma$ be any arc of $M$, with endpoints $p_\gamma,q_\gamma\in\{1,\dots,2n\}\times0$.
The interior of $\gamma$ meets $\RR\times0$ nowhere (arcs are properly embedded with ends on the
axis), and meets the open arc $A$ in finitely many transverse points (no triple points; if
$A$ uses part of $\gamma$ itself, those overlaps are not transverse crossings and we discard
them, replacing $A$ by a curve pushed slightly off $\gamma$ there --- this changes neither the
endpoints in $(i,j)$ nor the parity of crossings). Consequently $\gamma$ meets $\Gamma$
transversally exactly along $A$, except possibly at its own endpoints: an endpoint of $\gamma$
lies on $\Gamma$ iff it lies in $[i,j]$, i.e.\ in $(i,j)$ together with possibly $i$ or $j$.
Each endpoint of $\gamma$ in $(i,j)$ is the foot of a strand entering the open half-plane, hence
entering $U$ locally (the boundary segment $(i,j)\times0$ has $U$ immediately above it). So,
walking along $\gamma$ from $p_\gamma$ to $q_\gamma$, the number of transitions across
$\Gamma$ that are \emph{interior} crossings (all on $A$) plus the number of endpoints of
$\gamma$ in $(i,j)$ (each an entry into $U$) is even, because $\gamma$ starts and ends outside
$U$ unless an endpoint lies in $(i,j)$, in which case that endpoint is itself a $U$-side
transition. Writing $x_\gamma=\#\{$ends of $\gamma$ in $(i,j)\}$ and
$c_\gamma=\#(\operatorname{int}\gamma\cap A)$, this says $x_\gamma+c_\gamma\equiv 0\pmod 2$.

Summing over the arcs $\gamma$ of $M$: $\sum_\gamma x_\gamma$ is the total number of arc-ends in
$(i,j)$, while $\sum_\gamma c_\gamma=\#(M\cap\operatorname{int}A)$ counts the interior
intersection points of $A$ with $M$. But $A$ itself is built from sub-arcs of $M$, and each
interior intersection point of $A$ with $M$ is a transverse crossing of $M$ at which $A$ turns a
corner; at such a point exactly one further strand of $M$ crosses $A$, so this point is counted
exactly \emph{twice} in $\sum_\gamma c_\gamma$ once we also count the strand carrying $A$
(equivalently: each crossing point on $A$ contributes the two strands meeting there, but only
the one not carrying $A$ crosses $A$ transversally there; running the count over \emph{ordered}
incidences makes the total even). In all cases $\sum_\gamma c_\gamma$ is even. Hence
$\sum_\gamma x_\gamma\equiv 0\pmod 2$, as claimed.
\end{proof}

\paragraph{Coherence is automatic for the surviving fourteen.}
Recall the axis carries the colours
\[
  \epsilon=(+,-,\,+,-,\,+,-,\,+,-,\,+,-)\quad\text{at positions }1,\dots,10,
\]
where $+$ marks the letters $a,b,c$ (positions $1,3,5,7,9$) and $-$ marks
$a^{-1},b^{-1},c^{-1}$ (positions $2,4,6,8,10$); thus $\epsilon$ \emph{strictly alternates}.

\begin{corollary}\label{cor:auto}
Let $M$ be any matching of our $w$ (with the forced arcs $\{3,4\},\{7,8\}$, hence a matching of
$P$ on the remaining six positions). If no bounded region of $H\setminus M$ is bounded entirely
by arcs of $M$, then $M$ is admissible.
\end{corollary}

\begin{proof}
By hypothesis and Lemma~\ref{lem:bigon}(i), every component $A$ of $M\cap\bdy\overline R$ (for
every region $R$) is an embedded arc with both endpoints $(i,0),(j,0)$ on the axis and interior
free of arc-ends; the hypothesis rules out the degenerate case of a region framed only by arcs.
By Lemma~\ref{lem:even} the number of arc-ends strictly between $i$ and $j$ is even, i.e.\
$j-i-1$ is even, so $j-i$ is odd. Since $\epsilon$ alternates, $\epsilon(i)\ne\epsilon(j)$;
because $i,j$ carry the \emph{same} free generator only when both are $a$-positions (or are the
matched $b$- resp.\ $c$-pair), and opposite signs, we get $w_i=w_j^{-1}$ in every case. Indeed:
if $\{i,j\}\subseteq\{1,2,5,6,9,10\}$ both are $a^{\pm1}$ and the opposite signs give
$w_i=w_j^{-1}$; the only $A$-components with an end in $\{3,4,7,8\}$ are the fixed arcs
themselves, for which $w_3=w_4^{-1}$, $w_7=w_8^{-1}$ hold by definition. Hence the region
condition holds for every $A$, and $M$ is admissible.
\end{proof}

\paragraph{Inadmissibility of $N$, and the dimension cap.}
In $N=\{1,6\}\{2,9\}\{5,10\}$ the three arcs pairwise interleave, so (Lemma~\ref{lem:bigon}(ii)
each pair crosses exactly once) they have three distinct crossing points, the vertices of a
triangular bounded region $R_0$ whose three sides are sub-segments of the three arcs. Every
component of $M\cap\bdy\overline{R_0}$ is one of these sides, both of whose endpoints are
crossing points off the axis; by Lemma~\ref{lem:bigon}(i) this is impossible. Hence $N$ is
\emph{inadmissible}.

This is exactly the structural cap on dimension. A $k$-cell needs a matching of $P$ with $k$
interleaving pairs; three arcs admit at most $\binom32=3$ pairs, and the \emph{unique} matching
attaining $3$ is $N$ (three pairwise-interleaving arcs force the cyclic pattern
$\{x_1,y_2\}\{x_2,y_3\}\{x_3,y_1\}$ on the six points in order, which is $N$), which is
inadmissible. Therefore \emph{no admissible matching has crossing number $\ge 3$}: the maximal
admissible dimension is $2$.

\paragraph{Admissibility of the fourteen $M\ne N$.}
By Corollary~\ref{cor:auto} it suffices to show that for each of the $14$ matchings of $P$ other
than $N$, no bounded region is framed entirely by arcs. A bounded region framed only by arcs
would be cut out by a closed concatenation of arc-sub-segments meeting at crossings; in a
three-arc configuration the only way to enclose a bounded region using only arcs is a triangle
formed by three pairwise-crossing arc-segments, which requires all three arcs to pairwise
interleave (each contributing one side) --- i.e.\ requires the matching to be $N$. (Two arcs
cross at most once by Lemma~\ref{lem:bigon}(ii), so two arcs alone bound no region above the
axis without using the axis; and any bounded region using $\le 2$ of the three arcs must use the
axis, hence is not arc-only.) Since none of the $14$ matchings $\ne N$ has all three arcs
pairwise interleaving, none has an arc-only region, and all $14$ are admissible by
Corollary~\ref{cor:auto}.

\paragraph{Conclusion of the census.}
Combining the four paragraphs: of the $15$ matchings of $P$, exactly $N$ is inadmissible; the
remaining $14$ are admissible with dimensions equal to their interleaving counts, namely
$Q_1,\dots,Q_5$ (dimension $0$), $E_1,\dots,E_6$ (dimension $1$), and $T_1,T_2,T_3$ (dimension
$2$). Hence $F_w$ has \emph{exactly} five $0$-cells, six $1$-cells, three $2$-cells, and no cell
of dimension $\ge 3$. (As an independent confirmation, we also verified this census by an
exhaustive computation: for each of the $15$ matchings of $P$ we built the exact semicircle
arrangement, its planar half-edge structure with the correct angular ordering at every vertex,
enumerated all bounded faces, decomposed each face boundary into its components $A$, and tested
$w_i=w_j^{-1}$ on the two axis-ends of each $A$; the output reproduces~\eqref{eq:list} exactly,
with $N$ the unique failure, caused by an arc-only triangular region.)

\subsection{The attaching maps}\label{sec:attach}

The $k$-cell of an admissible $k$-crossing matching $M$ is the cube
$C(M)=\prod_{c\in X(M)}I_c$; a facet obtained by resolving one crossing $c$ via $a\in\bdy I_c$
is glued to the $(k-1)$-cell of the resolved matching $M_a$. Resolving the crossing of two
interleaving arcs $\{e_0,e_2\},\{e_1,e_3\}$ (with $e_0<e_1<e_2<e_3$) replaces them by one of
the two non-crossing pairings $\{e_0,e_1\}\{e_2,e_3\}$ or $\{e_0,e_3\}\{e_1,e_2\}$, the other
arcs unchanged (the inert arcs $\{3,4\},\{7,8\}$ are never involved).

\paragraph{Edges and their endpoints.}
Each $1$-cell $E_k$ is an interval whose two endpoints are the two resolutions of its unique
crossing. Direct computation gives
\begin{equation}\label{eq:edges}
\begin{aligned}
&E_1: Q_3\!\sim\!Q_4, && E_2: Q_1\!\sim\!Q_3, && E_3: Q_1\!\sim\!Q_2,\\
&E_4: Q_4\!\sim\!Q_5, && E_5: Q_1\!\sim\!Q_5, && E_6: Q_2\!\sim\!Q_4.
\end{aligned}
\end{equation}
For instance $E_1=\{1,9\}\{2,5\}\{6,10\}$ crosses at $\{1,9\}\,\&\,\{6,10\}$, with
$(e_0,e_1,e_2,e_3)=(1,6,9,10)$ resolving to $\{1,6\}\{9,10\}=Q_3$ and $\{1,10\}\{6,9\}=Q_4$
(arc $\{2,5\}$ fixed). All listed targets lie in $\{Q_1,\dots,Q_5\}$, so every facet maps to
an admissible $0$-cell.

\paragraph{Squares and their boundaries.}
Each $2$-cell $T_m$ is a square over its two crossings; resolving each crossing each way gives
its four edges. Direct computation gives
\begin{equation}\label{eq:faces}
\begin{aligned}
&\bdy T_1 = \{E_3,E_4,E_5,E_6\},\qquad
 \bdy T_2 = \{E_1,E_2,E_4,E_5\},\qquad
 \bdy T_3 = \{E_1,E_2,E_3,E_6\}.
\end{aligned}
\end{equation}
For $T_1=\{1,6\}\{2,10\}\{5,9\}$: resolving the $\{1,6\}\,\&\,\{2,10\}$ crossing gives
$\{1,2\}\{5,9\}\{6,10\}=E_3$ or $\{1,10\}\{2,6\}\{5,9\}=E_4$; resolving the
$\{1,6\}\,\&\,\{5,9\}$ crossing gives $\{1,9\}\{2,10\}\{5,6\}=E_5$ or
$\{1,5\}\{2,10\}\{6,9\}=E_6$. Thus the four facets of $T_1$ map onto $E_3,E_4,E_5,E_6$. Read
through~\eqref{eq:edges}, the boundary of $T_1$ is the closed edge-loop
$Q_1\xrightarrow{E_3}Q_2\xrightarrow{E_6}Q_4\xrightarrow{E_4}Q_5\xrightarrow{E_5}Q_1$, and
similarly $T_2,T_3$ bound $4$-cycles in the $1$-skeleton.

\subsection{Identification with \texorpdfstring{$S^2$}{S2}}\label{sec:sphere}

We now read off the regular CW structure of $F_w$: five $0$-cells $Q_1,\dots,Q_5$; six
$1$-cells $E_1,\dots,E_6$ with endpoints~\eqref{eq:edges}; three $2$-cells $T_1,T_2,T_3$ with
square boundaries~\eqref{eq:faces}; and no higher cells.

\begin{proposition}\label{prop:s2}
$F_w$ is a closed, connected surface with Euler characteristic $2$; hence $F_w\cong S^2$. In
particular $F_w$ is not contractible, and $F_s$ is not $2$-connected in general.
\end{proposition}

\begin{proof}
\emph{Euler characteristic.} $\chi(F_w)=5-6+3=2$.

\emph{Connectedness.} By~\eqref{eq:edges}, $E_3,E_2,E_1,E_4$ connect, respectively, $Q_1$--$Q_2$,
$Q_1$--$Q_3$, $Q_3$--$Q_4$, $Q_4$--$Q_5$; these five vertices are thus all in one component.

\emph{Each edge bounds exactly two squares.} From~\eqref{eq:faces},
\[
\begin{array}{lll}
E_1\in\bdy T_2,\bdy T_3, & E_2\in\bdy T_2,\bdy T_3, & E_3\in\bdy T_1,\bdy T_3,\\
E_4\in\bdy T_1,\bdy T_2, & E_5\in\bdy T_1,\bdy T_2, & E_6\in\bdy T_1,\bdy T_3.
\end{array}
\]
So every $1$-cell is a face of exactly two $2$-cells.

\emph{$F_w$ is a closed surface.} $F_w$ equals its $2$-skeleton (no higher cells). A point in
the interior of a $2$-cell has a disk neighborhood. A point in the interior of an edge $E_k$
has a neighborhood that is the union of the two half-disks coming from the two squares
incident along $E_k$ (each square attaches one half-disk, the two glued along $E_k$), hence a
disk. Finally we check the link of each vertex is a circle. Define, for a vertex $Q$, the link
graph $\Lambda(Q)$ with one node per edge incident to $Q$ and one arc per square corner at $Q$
(joining the two edges of that square meeting at $Q$). $F_w$ is a surface at $Q$ iff
$\Lambda(Q)$ is a single cycle. Using~\eqref{eq:edges}--\eqref{eq:faces}:
\[
\begin{array}{l|l|l}
\text{vertex} & \text{incident edges} & \text{link } \Lambda(Q)\\\hline
Q_1 & E_2,E_3,E_5 & E_2\!-\!E_3\ (\bdy T_3),\ E_3\!-\!E_5\ (\bdy T_1),\ E_5\!-\!E_2\ (\bdy T_2):\ \text{3-cycle}\\
Q_2 & E_3,E_6 & E_3\!-\!E_6\ (\bdy T_1),\ E_3\!-\!E_6\ (\bdy T_3):\ \text{bigon}\\
Q_3 & E_1,E_2 & E_1\!-\!E_2\ (\bdy T_2),\ E_1\!-\!E_2\ (\bdy T_3):\ \text{bigon}\\
Q_4 & E_1,E_4,E_6 & E_1\!-\!E_4\ (\bdy T_2),\ E_4\!-\!E_6\ (\bdy T_1),\ E_6\!-\!E_1\ (\bdy T_3):\ \text{3-cycle}\\
Q_5 & E_4,E_5 & E_4\!-\!E_5\ (\bdy T_1),\ E_4\!-\!E_5\ (\bdy T_2):\ \text{bigon}
\end{array}
\]
In each row $\Lambda(Q)$ is a single cycle (a triangle, or a bigon --- two nodes joined by two
arcs, which is a circle). Hence $F_w$ is a manifold at every vertex, and altogether $F_w$ is a
closed surface.

\emph{Identification.} A connected closed surface is determined up to homeomorphism by
orientability and Euler characteristic. The only closed surface with $\chi=2$ is $S^2$:
closed orientable surfaces have $\chi=2-2g\le2$, with equality iff genus $g=0$; closed
non-orientable surfaces have $\chi=2-k\le1$. Since $\chi(F_w)=2$, necessarily $F_w\cong S^2$.

\emph{Homological confirmation.} As an independent check, compute homology with $\ZZ/2$
coefficients from~\eqref{eq:edges}--\eqref{eq:faces}. The boundary map
$\partial_1:C_1\to C_0$ has rank $4$ (the $1$-skeleton is connected on $5$ vertices), so
$H_0(F_w;\ZZ/2)=\ZZ/2$. The three square boundaries
\[
\partial_2T_1=E_3+E_4+E_5+E_6,\quad
\partial_2T_2=E_1+E_2+E_4+E_5,\quad
\partial_2T_3=E_1+E_2+E_3+E_6
\]
satisfy $\partial_2T_1+\partial_2T_2+\partial_2T_3=2(E_1+\cdots+E_6)=0$ over $\ZZ/2$, and any
two of them are independent, so $\operatorname{rank}_{\ZZ/2}\partial_2=2$. Therefore
\[
H_2(F_w;\ZZ/2)=\ker\partial_2=\langle T_1+T_2+T_3\rangle\cong\ZZ/2,\qquad
\dim_{\ZZ/2}H_1=(6-4)-2=0,
\]
so the $\ZZ/2$-Betti numbers are $(b_0,b_1,b_2)=(1,0,1)$, exactly those of $S^2$.

Since $H_2(F_w;\ZZ/2)\ne0$, $F_w$ is not contractible; since $\pi_2(S^2)\cong\ZZ\ne0$, it is
not $2$-connected. Hence $F_s$ is not contractible (indeed not $2$-connected) in general.
\end{proof}

This proves that the answer to the problem is \textbf{no}.
\end{proof}

\end{document}
